The model draws on three literatures that have developed largely apart. We review each for what it establishes and what it leaves open, and we close by identifying the connection none of them makes.

\subsection{Generations at Work: What the Evidence Supports}

The sociological idea of a generation comes from \citet{mannheim1952problem}, who argued that people born in the same period share a location in historical time and may, if they are exposed to the same formative events, develop a shared outlook. Mannheim did not claim that birth cohorts automatically become generations in this richer sense; shared location was a precondition, not a guarantee. Management research has often skipped that qualification and treated cohort labels as if they carried stable traits.

The empirical record on generational differences at work is modest. A large time-lag study of American high school seniors found that the value placed on leisure rose across cohorts and work centrality fell, but it did not find that the most recent cohort it studied placed more weight on altruistic work values such as helping others or contributing to society \citep{twenge2010generational}. That finding is worth keeping in mind, because popular accounts of younger workers often assume the opposite. A meta-analysis of generational differences in job satisfaction, organizational commitment, and turnover intentions concluded that the differences were small or close to zero in most comparisons \citep{costanza2012generational}. Reviews of work values and of the wider workplace literature describe the evidence as mixed and largely descriptive, and both point to a recurring difficulty in separating the effect of cohort from the effects of age and period \citep{parry2011generational,lyons2014generational}.

A more critical line of work questions whether generations are a useful explanatory category for individual behavior at all. \citet{rudolph2017considering} proposed a lifespan developmental perspective as an alternative, arguing that the processes usually attributed to generation are better understood as developmental and contextual. \citet{rudolph2018leadership} reached a similar conclusion for leadership research, and \citet{rudolph2021generations} set out a series of common claims about generations that they argued are not supported by evidence. A different response to the same problem is to treat generation as an identity that becomes relevant in some organizational situations and not others \citep{joshi2010unpacking}, which shifts attention from what a cohort is like to when cohort membership matters. In our reading, research that links Generation Z specifically to workplace ethics is thinner still, and much of what circulates comes from practitioner surveys rather than peer-reviewed studies.

For our purposes, this literature establishes a caution rather than a finding. It gives little support for building a model on the premise that Generation Z holds stronger ethical values than earlier cohorts. It leaves open a narrower question that it has rarely asked: whether cohorts differ in the channels through which they act on ethical concerns, and whether any such difference survives once age and career stage are taken into account.

\subsection{Ideological Psychological Contracts and Responses to Breach}

Psychological contract research began with the insight that employees hold beliefs about mutual obligations that go beyond any formal agreement \citep{rousseau1989psychological}. Later work traced how these beliefs form, locating their origins partly in pre-entry schemas and in promises communicated through the recruitment and socialization process \citep{rousseau2001schema}. A central distinction separates breach, the perception that an obligation has not been met, from violation, the emotional response of anger and betrayal \citep{morrison1997employees}. Longitudinal evidence showed that breach was more likely when newcomers had little interaction with organizational agents before being hired, and that breach turned into more intense violation when it was attributed to deliberate reneging and handled unfairly \citep{robinson2000development}. A meta-analysis confirmed that breach is associated with a wide range of attitudes and behaviors and identified violation and mistrust as affective links in that chain \citep{zhao2007impact}.

\citet{thompson2003violations} widened the theory by proposing ideological currency: commitments by the organization to a cause or principle that the employee values and that benefits parties beyond the employee. This extension brought ethical commitments inside the psychological contract. A recent review lists ideological currency among the directions in which the field needs further research \citep{coyleshapiro2019psychological}. Empirical work on how employees respond to ideological breach has begun to accumulate. One study found that breach prompted pro-social rule breaking, which the authors interpret as corrective behavior on behalf of the ideology \citep{yang2022going}. Another, in a hospital setting, found that ideological breach led employees to ruminate and then to reaffirm their own values through compensatory effort on behalf of the cause \citep{deng2023serving}.

These studies examine responses directed at the work itself or at the cause. Earlier research on non-ideological contracts looked at responses directed at the organization: managers who experienced violations reported more exit, voice, and neglect and less loyalty, and the availability of attractive job alternatives moderated the link to exit \citep{turnley1999impact}. What this literature does not examine is the channel through which an employee who wants to act on a broken ethical commitment chooses to do so, nor whether that choice depends on how the commitment was communicated in the first place. The formation research tells us that promises are conveyed through many channels, and the breach research tells us that how employees interpret a breach shapes the emotional response, yet the two have not been joined to ask whether the public or private origin of an ethical promise changes what employees do when it is broken.

\subsection{Ethical Voice, Silence, and the Choice of Channel}

Behavioral ethics research asks how people in organizations come to act rightly or wrongly \citep{trevino2006behavioral}. Two streams within and alongside it address what employees do when they see something wrong. The voice stream studies upward communication of concerns and suggestions and its opposite, silence \citep{morrison2014employee,morrison2023employee}. It has shown that silence can be widespread and collectively reinforced \citep{morrison2000organizational} and that the decision to speak rests on judgments about safety and efficacy. The whistleblowing stream studies disclosure of wrongdoing to parties who can act on it, and its standard definition includes disclosures both inside and outside the organization \citep{near1985organizational}. A book-length synthesis summarized decades of this research \citep{miceli2008whistleblowing}, and a review organized the field around who blows the whistle, about what, how, why, and to whom \citep{culiberg2017evolution}.

The question of channel appears most directly in the whistleblowing stream. A comparison of dismissed whistleblowers found that those who went outside the organization had shorter tenure and stronger evidence than those who reported internally, were more effective in changing practices, and suffered more extensive retaliation \citep{dworkin1998internal}. A study of intended responses to observed wrongdoing found that stronger ethical cultures were associated with less intended inaction and less intended external whistleblowing, and with more intended internal reporting \citep{kaptein2011inaction}. A meta-analysis found that personal, situational, and wrongdoing characteristics predicted whistleblowing intentions more strongly than actual whistleblowing \citep{mesmermagnus2005whistleblowing}, a reminder that stated intentions are a weak guide to behavior.

Voice research offers a further distinction that bears on ethical concerns. Voice can promote a new idea or it can object to something harmful that is already happening, and the second kind, prohibitive voice, carries more interpersonal risk because it implies criticism of current practice or of the people responsible for it. In a two-wave study, psychological safety was more strongly related to prohibitive than to promotive voice \citep{liang2012psychological}. Raising an ethical objection to an employer's conduct is a clear case of prohibitive voice. The same research tradition also shows that silence is not a single state: employees may withhold concerns because they have given up, because they fear the consequences, or because they wish to protect others \citep{vandyne2003conceptualizing}. These distinctions matter for the model because they imply that the internal channel is least attractive precisely for the kind of concern at issue here.

Social media have added a new external channel. Commentators have pointed out that employees can now reach large audiences with their views of the employer, for better or worse \citep{miles2014employee}, and a study in business ethics examined how an organization's social media strategy and policy affect employees' sensitivity to the risks of whistleblowing through social media \citep{xiao2022employee}. Communication research explains part of the channel's reach: social media make content visible to wide audiences and keep it available long after it is posted \citep{treem2013social}. This work treats the social media channel as a risk to be managed. It does not ask what makes the channel attractive to a given employee after a given breach.

\medskip

Across the three literatures, the missing link is the relationship between where an ethical promise is made and where its breach is contested. Generational research has focused on value levels and has not modeled channel choice. Ideological contract research has studied how employees feel and what they do for the cause after breach, but has not examined the venue of their response or the venue of the original promise. Voice and whistleblowing research has distinguished internal from external channels and has identified ethical culture and safety as influences on the choice, but has not asked whether the promise's origin matters. The nearest existing work combines ideological contract and signaling reasoning to study employees' online whistleblowing about their employers' public stances on social issues, with the practical aim of guiding internal communication about those stances \citep{yang2026managing}. It shares our interest in public promises and public responses. It does not model internal voice safety as a boundary, exit and silence as alternative outlets, or the rival cohort and career-stage explanations that the Generation Z question requires. Those are the elements our model adds.
